\section{Discussion}\label{sec:discussion}

The screen is a first stage. Its intended
pipeline reduces $p$ coordinates to a manageable $d$, after which a joint
method --- the central tail-index subspace of \citet{gardesPodgorny2025}, a
structured tail regression \citep{wangTsai2009,sasakiTaoWang2026}, or
multivariate local estimation --- can resolve interactions and correlated
neighbours among the retained variables. The two stages solve different
problems: rapid elimination of evidently irrelevant coordinates scales
to large $p$, while accurate joint tail inference does not.

The scope of a marginal screen for this target is delimited exactly
rather than assumed. An active coordinate every one of whose fibres
meets the global argmax set leaves no trace in the score, at any sample
size; Proposition~\ref{prop:geometry} characterizes this case exactly,
in terms of the argmax projections, where a signal-strength condition
would only have hidden it. Knowing the boundary is what allows a user
to tell a screen that has failed from a target that a marginal screen
cannot see, and it marks where iterative or grouped extensions, with
their own identification arguments, would be required. Two further
features of the target are worth stating in the same spirit. Projected
mixture bias is logarithmic, so (E2) covers every gap with
$\Delta_{\min}\log(1/\alpha)\to\infty$ and in particular every fixed
gap; whether gaps of the excluded order $1/\log n$ are detectable by
any estimator is a question about the target, not about this screen.
And the envelope near its maximum is estimated from observations whose
remaining active coordinates are simultaneously near their maximizing
values, so the sample cost grows with the number of active coordinates:
the weak-signal model A2 in Section~\ref{sec:simulation} measures it,
and aggregation over neighbouring tunings recovers most of it.

Appendix~\ref{app:primitive} derives the estimation conditions from
primitive assumptions that are verifiable on concrete models; what
remains is to establish whether the logarithmic projected bias, which
restricts the covered gaps to $\Delta_{\min}\log(1/\alpha)\to\infty$, is a
property of the target or of this estimator. The score functional also
invites variants: a minimum or a low quantile of the profile would
emphasize depressions confined to short intervals, at the price of new
uniformity arguments for empirical profile extrema.

In practice, tuning matters more than the theory prescribes. The tail
fraction trades Hill bias against variance; the bandwidth additionally
determines how often near-maximizing configurations of the other active
coordinates enter a window. We recommend the workflow of
Sections~\ref{sec:simulation} and~\ref{sec:realdata}: run the screen on
the block of nine neighbouring tunings and rank by the aggregated
minimum rank. The aggregation is deliberately conservative --- a
predictor is retained when at least one reasonable tuning finds
evidence for it --- which is the right direction of error for a
first-stage sure screen: false exclusion is final, while false
retention is corrected by the second-stage multivariate analysis. The
per-setting ranks describe tuning sensitivity and should accompany the
ranking as a diagnostic, not act as a second screening threshold.
Dependence deserves caution of its own: it can help, by concentrating
active configurations, or hurt, by promoting inactive neighbours, and
the direction is design-specific.

Projection turns a tail-index surface into an upper envelope. Variation
of that envelope is a scalable marginal signal for tail-index activity,
and its absence, characterized by the argmax projections, tells the
analyst exactly what a marginal screen can and cannot see.
